\subsection*{S19. Recursive and mutually recursive shapes}

A corpus whose structure refers to itself: trees, comment threads, file
systems, organisation charts.

\begin{lstlisting}[style=json]
{ "value": 1,
  "children": [ { "value": 2, "children": [] },
                { "value": 3, "children": [ { "value": 4,
                                             "children": [] } ] } ] }
\end{lstlisting}

\paragraph{Doing nothing is not rejection.}
JSON documents are finite trees, so any corpus has bounded depth and naive
inference emits one table per level observed --- \texttt{node},
\texttt{node\_children}, \texttt{node\_children\_children}. Nothing fails. But
the number of tables then depends on how deep the sample happened to go, and a
document arriving later that is deeper than the sample has nowhere to be put:
shredding becomes partial. This is the same defect that ruled out treating
maps as records in S9.

\paragraph{Self-recursion and mutual recursion differ in cost.}
A single self-referential table needs no special machinery: \texttt{val
children : t -> t list} refers to \texttt{t} within its own signature, which is
ordinary OCaml. Recursion running through two or more tables requires
\texttt{module rec}, which carries real restrictions.

\options
\begin{enumerate}[label=\textbf{(\alph*)}]
  \item \textbf{Detect and reject.} Note this still requires the detection
        machinery of (b); it merely errors rather than using it.
  \item \textbf{Detect and emit a recursive schema.}
  \item \textbf{No detection} --- one table per observed depth, as above.
  \item \textbf{Explicit annotation}, as in S9: the configuration names the
        recursive paths. \selected
\end{enumerate}

\decision{Explicit annotation. The configuration file identifies a path with
an ancestor path it folds into; inference then unifies the two rather than
generating a fresh table per level.}

\begin{lstlisting}[style=json]
recursive = [ { path = ".children[]", folds_into = "." } ]
\end{lstlisting}

The annotation identifies a \emph{pair} of paths, unlike S9's marking which
names a single path. Naming only the recursive path would require inferring
which ancestor it folds into, reintroducing the detection this option exists
to avoid.

The two paths unify into one table with a self-referencing foreign key. Root
rows have no parent, so the parent reference and the index are both partial:

\begin{center}
\begin{tabular}{@{}llll@{}}
\multicolumn{4}{@{}l}{\textbf{Table node}} \\
\hline
\texttt{id} & \texttt{parent\_id} & \texttt{idx} & \texttt{value} \\
\hline
1 & --- & --- & 1 \\
2 & 1 & 0 & 2 \\
3 & 1 & 1 & 3 \\
4 & 3 & 0 & 4 \\
\hline
\end{tabular}
\end{center}

\begin{lstlisting}[style=ml]
module Node : sig
  type id
  type t = { id : id; parent_id : id option; idx : int option;
             value : int }
  val get : db -> id -> t
  val children : db -> t -> t list
end
\end{lstlisting}

Where the recursion runs through more than one table, the generated modules
become mutually recursive:

\begin{lstlisting}[style=ml]
module rec Post : sig
  type id
  type t = { id : id; title : string }
  val get : db -> id -> t
  val comments : db -> t -> Comment.t list
end
and Comment : sig
  type id
  type t = { id : id; post_id : Post.id; idx : int; body : string }
  val get : db -> id -> t
  val post : db -> t -> Post.t
end
\end{lstlisting}

\paragraph{The annotation is validated, not trusted.}
After inference, the two paths named must actually have compatible shapes. If
they do not, the annotation is an error naming both paths and the mismatch.
Otherwise a mistaken annotation would produce a schema the data does not fit.

\rationale{Detection is sample-dependent, which is what rules out (b) and (a).
Whether \texttt{.children[]} has the same inferred shape as the root depends on
what was observed: if no child in the sample has grandchildren, then
\texttt{.children[].children} was only ever empty, its element type is unknown
under S15, and it need not compare equal to the root's shape. The schema would
then be recursive or not according to how deep the corpus happened to go ---
the data-dependence problem of S9.

Annotation avoids this the same way S9 does, and reuses the configuration
mechanism already introduced there. It also keeps shredding total: once a path
is known to be recursive, a document deeper than anything in the corpus still
has a table to go in.}

\limitation{As in S9, this gives determinism rather than protection. Recursive
data that is not annotated still lowers as one table per observed depth, and
is still partial on deeper documents; the user only finds out by inspecting
the generated modules. A detector could be added later as an advisory hint,
subject to the same constraint as S9's: it may suggest an annotation for a
person to accept, and must never change what is generated on its own.}
